\section{Evaluation protocol}
\label{sec:protocol}

Fish has a single chance node --- the deal --- followed by a public
transcript, and that structure permits an unusually strong variance-reduction
design. This section fixes the experimental unit, the interval construction,
the fitting/evaluation separation, and the metrics of \S\ref{sec:results}.

\subsection{Six-rotation duplicate blocks}
\label{sec:protocol-blocks}

Seats alternate between the two teams, so the cyclic group of six seat
rotations of a fixed deal forms a balanced block. The three odd rotations
exchange which team holds which hand-triple; the three even rotations permute
hands within a team. Playing all six rotations of every deal therefore has
each policy hold every hand-triple of that deal exactly once, and the
intrinsic advantage of the deal cancels from the comparison rather than merely
being averaged down. A matchup reported over $D$ deals is $6D$ games organised
into $D$ blocks of six.

This is a stronger control than the two-orientation swap used in the v0.3
study \cite{fishbot03}, which balances the team label only and leaves the
assignment of individual hands within a team uncontrolled. The frozen-policy
ablations of \S\ref{sec:ablations} use two rotations rather than six, for
budget reasons; that is stated in the caption of Table~\ref{tab:ablations} and
is a weaker block than the primary bank uses.

\subsection{Clustered interval construction}
\label{sec:protocol-ci}

The six rotations of one deal are one correlated cluster: they share the same
card distribution and differ only in seat assignment. Resampling \emph{games}
would treat those six observations as independent and understate the variance.
Every interval reported in this paper therefore resamples \emph{deals}: a cluster
bootstrap of the per-deal win count for a single arm, and a paired bootstrap of
the per-deal difference for a comparison between two configurations played on the
same deals, both implemented in \filepath{engine/src/arena.hpp} and both taking
the $2.5\%$ and $97.5\%$ quantiles of \numBootstrapDraws{} resamples.

They are therefore \emph{percentile} bootstrap intervals, with the deal as the
resampling cluster, and they inherit the known coverage error of that
construction: \cite{jordan-rleval} measures raw percentile intervals failing at
$5.7$--$11.2\%$ against a nominal $5\%$ rate. We keep the percentile
construction for two reasons. First, clustering by deal is the correction that
matters most in this design, because the within-deal correlation across the six
rotations is large and ignoring it would understate the variance by a
substantial factor, whereas the percentile-versus-BCa distinction is a
second-order adjustment at the sample sizes used here. Second, we have not
implemented the bias-corrected and accelerated (BCa) or bootstrap-$t$
corrections, so no claim of nominal coverage is made: the intervals should be
read as approximate, and differences whose intervals sit close to zero should not
be treated as resolved. Wilson intervals on the raw game counts are computed
alongside the bootstrap for continuity with the v0.3 paper; they ignore
clustering and are reported only as a reference point.

\subsection{Seed separation and what was held out}
\label{sec:protocol-seeds}

Fitting used recorded base seeds $20260821$ (round~1), $770077$ (round~2),
$313131$ (round~3) and $888111$ (round~4). The base seed of round~5, from
whose trace the shipping configuration was drawn, is not recorded in any
committed script or artifact; this is a reproducibility gap and is documented
as one in \S\ref{app:repro}. Selection over the last six generation means of
the round-5 trace used a further validation bank, seed $1357911$.

Every number reported in \S\ref{sec:results} comes from a bank disjoint from
all of those: $90210$ for the head-to-head, the declaration pre-gate audit and
the arbitration sensitivity run; $515151$ for the round-robin matrix; $606060$
for the frozen-policy ablations; $717171$ for calibration; $828282$, $838383$
and $848484$ for the rule dialects; $20260820$ for the v0.3 port validation;
$959595$ for the belief substitution under an unfitted policy; $464646$ for
action-cap incidence; $31415$ for the value-function fit; $20260822$ for the
brute-force oracle; and $515253$ (response fitting) with $6543210$ (response
evaluation) for the local response probe. The configuration was frozen before
any evaluation bank was run.

That separation is narrower than ``held out'' usually implies: deals were held
out --- no deal used in any reported experiment appears in any fitting or
selection bank --- but most opponent identities were not, and
\S\ref{sec:results-h2h} states the consequence for the evaluation panel and
scopes every result of \S\ref{sec:results} accordingly.

\subsection{Metrics}
\label{sec:protocol-metrics}

Win rate is the primary metric, but it is not sufficient on its own: a policy
can win while converting a smaller fraction of its asks, and
\S\ref{sec:results-h2h} reports a matchup in which exactly that happens. For
every matchup we therefore also report \emph{mean half-suits} won out of nine,
giving the margin and not only the sign of the result; \emph{ask accuracy},
the fraction of asks that transfer a card; and \emph{declaration accuracy and
rate}, declarations per game and the fraction scored correctly. The engine
tallies voluntary and forced endgame declarations separately; the win-rate
tables report the voluntary counters only, the calibration study pools the
two, and forced declarations run at well under a tenth per game on the primary
bank, so that contamination is small but is not zero. \emph{Out-of-turn
declarations} are reported per game; they are impossible under the v0.3
dialect, which permits declarations only on turn.

\emph{Forecast calibration} is the Brier score, log loss, expected calibration
error and a decile reliability table for the policy's own $\Pr[\text{ask
succeeds}]$ and $\Pr[\text{allocation correct}]$ forecasts, descriptive
properties of the estimators computed as point estimates without clustered
uncertainty. The \emph{worst-case profile} is the minimum win rate over the
opponent panel, not only the mean: a policy that averages well and collapses
against one playstyle is not well described by its mean. The \emph{local
response probe} is a policy fitted specifically against a frozen target,
inside the same policy class and with the same optimiser and budget; the score
such a response achieves is an achieved \emph{lower} bound on the
exploitability of the frozen target within the searched policy class, since a
larger budget, a better optimiser or a policy class outside the one searched
could all achieve more. It is a restricted-policy-class diagnostic and not an
equilibrium computation.
